% ---- pg-526 （书页 514） ----

\noindent\hfill{\zihao{6}续表}\par\vspace{2pt}

\begin{mvtable}
\mvfit{%
\begin{tabular}{cccccccccccc}
\toprule
\multicolumn{4}{c}{$\pi_1$} & \multicolumn{4}{c}{$\pi_2$} & \multicolumn{4}{c}{$\pi_3$}\\
\cmidrule(lr){1-4}\cmidrule(lr){5-8}\cmidrule(lr){9-12}
萼片长 & 萼片宽 & 花瓣长 & 花瓣宽 & 萼片长 & 萼片宽 & 花瓣长 & 花瓣宽 & 萼片长 & 萼片宽 & 花瓣长 & 花瓣宽\\
度 $x_1$ & 度 $x_2$ & 度 $x_3$ & 度 $x_4$ & 度 $x_1$ & 度 $x_2$ & 度 $x_3$ & 度 $x_4$ & 度 $x_1$ & 度 $x_2$ & 度 $x_3$ & 度 $x_4$\\
\midrule
5.0 & 3.5 & 1.6 & 0.6 & 5.0 & 2.3 & 3.3 & 1.0 & 6.8 & 3.2 & 5.9 & 2.3\\
5.1 & 3.8 & 1.9 & 0.4 & 5.6 & 2.7 & 4.2 & 1.3 & 6.7 & 3.3 & 5.7 & 2.5\\
4.8 & 3.0 & 1.4 & 0.3 & 5.7 & 3.0 & 4.2 & 1.2 & 6.7 & 3.0 & 5.2 & 2.3\\
5.1 & 3.8 & 1.6 & 0.2 & 5.7 & 2.9 & 4.2 & 1.3 & 6.3 & 2.5 & 5.0 & 1.9\\
4.6 & 3.2 & 1.4 & 0.2 & 6.2 & 2.9 & 4.3 & 1.3 & 6.5 & 3.0 & 5.2 & 2.0\\
5.3 & 3.7 & 1.5 & 0.2 & 5.1 & 2.5 & 3.0 & 1.1 & 6.2 & 3.4 & 5.4 & 2.3\\
5.0 & 3.3 & 1.4 & 0.2 & 5.7 & 2.8 & 4.1 & 1.3 & 5.9 & 3.0 & 5.1 & 1.8\\
\bottomrule
\end{tabular}
}
\end{mvtable}

\noindent{\zihao{6}资料来源: Anderson[1].}

\noindent (a)\quad 在 $(x_2,x_4)$ 变量空间对数据作图. 这三组观测数据是否表现为二元正态分布?

\noindent (b)\quad 假定样本来自有相同协方差矩阵的二元正态总体. 检验假设 $H_0$: $\boldsymbol{\mu}_1=\boldsymbol{\mu}_2=\boldsymbol{\mu}_3$, $H_1$: 至少有一个 $\boldsymbol{\mu}_i$ 与其他均值不同, 显著性水平取 $\alpha=0.05$. 在这种情况下, 关于公共协方差矩阵的假定是否合理? 请解释.

\noindent (c)\quad 假定总体服从二元正态分布, 构造由\hyperref[eq:11-47]{式(11-47)}给出的二次判别得分 $\hat{d}_i^Q(\mathbf{x})$, 假定 $p_1=p_2=p_3=1/3$. 用法则(11-48)将 $\mathbf{x}_0'=[3.5,1.75]$ 分类到总体 $\pi_1,\pi_2$ 或 $\pi_3$.

\noindent (d)\quad 假定协方差矩阵 $\boldsymbol{\Sigma}_i$ 对三个二元正态总体均相同. 构造由\hyperref[eq:11-51]{式(11-51)}所给出的线性判别得分 $\hat{d}_i(\mathbf{x})$, 并根据\hyperref[eq:11-52]{式(11-52)}将 $\mathbf{x}_0'=[3.5,1.75]$ 分类到 $\pi_1,\pi_2$ 或 $\pi_3$. 取 $p_1=p_2=p_3=1/3$. 将结果与(c)中结果比较. 你倾向于哪种方法? 请解释.

\noindent (e)\quad 假定三个总体均为二元正态总体, 有相同的协方差阵, 且 $p_1=p_2=p_3=1/3$. 用法则(11-56)将 $\mathbf{x}_0'=[3.5,1.75]$ 分配到 $\pi_1,\pi_2$ 或 $\pi_3$. 将结果与(d)的结果比较. 将由法则(11-56)中的 $\hat{d}_{ki}(\mathbf{x}_0)$ 所决定的分类域 $\hat{R}_1,\hat{R}_2$ 和 $\hat{R}_3$ 画在(a)中所作的图上.

\noindent (f)\quad 用(d)中得出的线性判别函数将样本观测值分类. 计算 APER 及 $\hat{E}(\mathrm{AER})$. (为计算后者, 你需要用拉亨布鲁克的提留方法[见\hyperref[eq:11-57]{式(11-57)}]).

\noindent 11.28\quad 达罗赫和摩西曼(见文献[6])认为, \hyperref[tab:11.5]{表 11.5} 中所列的三种鸢尾可以只根据“形状”或与标度无关的信息来判别. 令 $Y_1=X_1/X_2$ 为萼片形状, 而 $Y_2=X_3/X_4$ 为花瓣形状.

\noindent (a)\quad 将数据在 $(\log Y_1,\log Y_2)$ 变量空间作散布图. 三组观测数据是否表现为二元正态分布?

\noindent (b)\quad 假定总体服从二元正态分布具有相同的协方差矩阵, 再假定 $p_1=p_2=p_3=1/3$, 构造由\hyperref[eq:11-51]{式(11-51)}所给的线性判别得分 $\hat{d}_i(\mathbf{x})$. 分三种情况: 同时用 $\log Y_1$ 和 $\log Y_2$, 分别用其中之一. 计算 APER.

\noindent (c)\quad 用(b)中的线性判别函数计算期望 AER 的提留估计值, 将结果填入右表.

\begin{mvtable}
\mvfit{%
\begin{tabular}{cc}
\toprule
变量 & 错分率\\
\midrule
$\log Y_1$ & \\
$\log Y_2$ & \\
$\log Y_1,\log Y_2$ & \\
\bottomrule
\end{tabular}
}
\end{mvtable}

将上面的错分率与例 11.12 中得到的结果进行比较. 是否关于形状的信息可以单独作为区分这三种鸢尾
